\chapter{Experiments}
\label{experiments}

\minitoc{}

\section{Introduction}

Experimental validation of theorical work is a crucial step in the development of new technologies
and methods. This chapter presents the experiments conducted to validate the theoretical work
presented in \cref{cable_modelling,robot_chain_control}.

The first experimental campaign presented in
\cref{e_augmented_catenary_model_experimental_validation} was conducted to gather the dataset used
in \cref{cm_acm_experimental_validation} in order to quantify the accuracy of the proposed model in
comparaison to actual cables in underwater conditions. In addition to using a varied set of cables
changes experimental conditions such as the speed of the robot and the direction of motion.

The second experimental campaign presented in \cref{e_robot_chain_control_experiment} was conducted
to validate the feasability of the proposed control methods and simulations presented in
\cref{robot_chain_control}. The robot chain used was scaled down from the three Bluerovs with a
surface vehicle to a pair of Bluerovs for practical reasons.

\section{Augmented catenary model experimental validation}
\label{e_augmented_catenary_model_experimental_validation}

The motivation and objective of these experiments were to study the accuracy of catenary models in
estimating the shape of tethers in water under dynamic conditions. This model is required for
model-based control of tethered underwater vehicles, such as obstacle avoidance.

\subsection{Experimental setup}
\label{e_acmev_experimental_setup}

The experiment was conducted in a \qty{7.2}{\meter} long, \qty{4.2}{\meter} wide, and
\qty{3}{\meter} deep tank (\cref{e_acmev_es_robot_in_tank}). The cable's and robot's positions were
recorded by a 5-camera, \qty{100}{\Hz} Qualisys\footnote{\url{www.qualisys.com/cameras/underwater}}
underwater motion capture system. The calibration of this system provided a standard deviation of
\qty{2.2}{\mm} across the tank's calibrated volume for a known object's length. The robot and each
cable were equipped with passive reflective markers
(\cref{e_acmev_es_robot_markers,e_acmev_es_all_cables}), and the cable's markers were evenly spaced
out by \qty{20}{\centi\meter}. The cables were attached to the robot and to a fixed point on the
side of the tank with thin wire
(\cref{e_acmev_es_robot_cable_attachment_point,e_acmev_es_rod_cable_attachment_point}) so that the
joint between the tether and its attachment point closely approximated a theoretical spherical
joint. The cables used are shown in \cref{e_acmev_es_all_cables}, and their characteristics are
listed in \cref{e_acmev_es_cables_characteristics}.

\begin{figure}
	\centering
	\begin{minipage}{.49\figurewidth}
		\centering
		\includegraphics[width=.9\textwidth]{catenary_experiment_bluerov_picture_2.jpg}
		\captionof{figure}{Robot equipped with passive reflective markers.}
		\label{e_acmev_es_robot_markers}
	\end{minipage}
	\begin{minipage}{.49\figurewidth}
		\centering
		\includegraphics[width=.9\textwidth]{catenary_experiment_bluerov_picture_1.jpg}
		\captionof{figure}{Robot in the tank.}
		\label{e_acmev_es_robot_in_tank}
	\end{minipage}
	\marginnote{Brice performed admirably during the two weeks of experiments.}
\end{figure}

\begin{figure}
	\centering
	\begin{minipage}{.49\figurewidth}
		\centering
		\includegraphics[width=.9\textwidth]{catenary_experiment_robot_cable_attachment_point.jpg}
		\captionof{figure}{Robot cable attachment.}
		\label{e_acmev_es_robot_cable_attachment_point}
	\end{minipage}
	\begin{minipage}{.49\figurewidth}
		\centering
		\includegraphics[width=.9\textwidth]{catenary_experiment_rod_cable_attachment_point.jpg}
		\captionof{figure}{Fixed point cable attachment.}
		\label{e_acmev_es_rod_cable_attachment_point}
	\end{minipage}
\end{figure}

The experiment was repeated three times for each cable under six experimental conditions, enabling
the testing of more or less taut cables with accelerations of different amplitudes. Each
experimental condition was a combination of a direction of motion (sway or surge), a speed
(\qty{0.3}{\m\per\s} or \qty{0.6}{\m\per\s}), and a starting point characterized by its horizontal
distance from the fixation point (approximately \qty{1.5}{\m} or \qty{2}{\m}).

In every experiment, the robot received the following commands:
\begin{enumerate}
	\item stabilization with auto-depth at \qty{1}{\m} depth for \qty{2}{\s},
	\item abrupt start with a constant surge or sway command applied for \qty{2.5}{\s},
	\item reversing of the thrusters for \qty{3}{\s} with a doubled velocity in the opposite direction to
	      achieve an abrupt stop,
	\item slowing down the thrusters to half the initial velocity for \qty{0.2}{\s}, and
	\item deactivation of the auto-depth control and complete stop of the thrusters after \qty{13}{\s}. The
	      robot then drifted slowly.
\end{enumerate}

\begin{figure}
	\renewcommand\thesubfigure{\arabic{subfigure}}
	\centering
	\begin{subfigure}{.48\figurewidth}
		\centering
		\begin{subfigure}{.48\textwidth}
			\centering
			\includegraphics[width=\textwidth]{cable_coaxial.jpg}
		\end{subfigure}
		\begin{subfigure}{.48\textwidth}
			\centering
			\includegraphics[width=\textwidth]{cable_coaxial_top.jpg}
		\end{subfigure}
		\caption{}
		\label{e_acmev_es_coaxial}
	\end{subfigure}
	\begin{subfigure}{.48\figurewidth}
		\centering
		\begin{subfigure}{.48\textwidth}
			\centering
			\includegraphics[width=\textwidth]{cable_four_pairs.jpg}
		\end{subfigure}
		\begin{subfigure}{.48\textwidth}
			\centering
			\includegraphics[width=\textwidth]{cable_four_pairs_top.jpg}
		\end{subfigure}
		\caption{}
		\label{e_acmev_es_four_pairs}
	\end{subfigure}
	\begin{subfigure}{.48\figurewidth}
		\centering
		\begin{subfigure}{.48\textwidth}
			\centering
			\includegraphics[width=\textwidth]{cable_two_pairs.jpg}
		\end{subfigure}
		\begin{subfigure}{.48\textwidth}
			\centering
			\includegraphics[width=\textwidth]{cable_two_pairs_top.jpg}
		\end{subfigure}
		\caption{}
		\label{e_acmev_es_two_pairs}
	\end{subfigure}
	\begin{subfigure}{.48\figurewidth}
		\centering
		\begin{subfigure}{.48\textwidth}
			\centering
			\includegraphics[width=\textwidth]{cable_floating.jpg}
		\end{subfigure}
		\begin{subfigure}{.48\textwidth}
			\centering
			\includegraphics[width=\textwidth]{cable_floating_top.jpg}
		\end{subfigure}
		\caption{}
		\label{e_acmev_es_floating}
	\end{subfigure}
	\begin{subfigure}{.48 \figurewidth}
		\centering
		\begin{subfigure}{.48\textwidth}
			\centering
			\includegraphics[width=\textwidth]{cable_rope.jpg}
		\end{subfigure}
		\begin{subfigure}{.48\textwidth}
			\centering
			\includegraphics[width=\textwidth]{cable_rope_top.jpg}
		\end{subfigure}
		\caption{}
		\label{e_acmev_es_rope}
	\end{subfigure}
	\begin{subfigure}{.48\figurewidth}
		\centering
		\begin{subfigure}{.48\textwidth}
			\centering
			\includegraphics[width=\textwidth]{cable_weighted.jpg}
		\end{subfigure}
		\begin{subfigure}{.48\textwidth}
			\centering
			\includegraphics[width=\textwidth]{cable_weighted_top.jpg}
		\end{subfigure}
		\caption{}
		\label{e_acmev_es_weighted}
	\end{subfigure}
	\begin{subfigure}{.48\figurewidth}
		\centering
		\begin{subfigure}{.48\textwidth}
			\centering
			\includegraphics[width=\textwidth]{cable_chain.jpg}
		\end{subfigure}
		\begin{subfigure}{.48\textwidth}
			\centering
			\includegraphics[width=\textwidth]{cable_chain_top.jpg}
		\end{subfigure}
		\caption{}
		\label{e_acmev_es_chain}
	\end{subfigure}
	\begin{subfigure}{.48\figurewidth}
		\centering
		\begin{subfigure}{.48\textwidth}
			\centering
			\includegraphics[width=\textwidth]{cable_elastic.jpg}
		\end{subfigure}
		\begin{subfigure}{.48\textwidth}
			\centering
			\includegraphics[width=\textwidth]{cable_elastic_top.jpg}
		\end{subfigure}
		\caption{}
		\label{e_acmev_es_elastic}
	\end{subfigure}
	\caption{
		Cables:
		(\subref{e_acmev_es_coaxial}) coaxial cable,
		(\subref{e_acmev_es_four_pairs}) four-pair network cable,
		(\subref{e_acmev_es_two_pairs}) two-pair network cable,
		(\subref{e_acmev_es_floating}) floating braided rope,
		(\subref{e_acmev_es_rope}) sail rope,
		(\subref{e_acmev_es_weighted}) leaded seamstress rope,
		(\subref{e_acmev_es_chain}) plain steel chain,
		(\subref{e_acmev_es_elastic}) elastic.
	}
	\label{e_acmev_es_all_cables}
	\renewcommand\thesubfigure{\alpha{subfigure}}
\end{figure}

\begin{table}
	\centering
	\robustify\bfseries
	\sisetup{ % pour changer le nombre de chiffres après la virgule affichés après compilation
		round-mode = places,
		round-precision = 3}
	\caption{Cables characteristics}
	\begin{tabular}{@{}l|SSSS@{}}
		\toprule
		cable \#                    & \subref{e_acmev_es_coaxial} & \subref{e_acmev_es_four_pairs} & \subref{e_acmev_es_two_pairs} & \subref{e_acmev_es_floating} \\ \midrule
		diameter / \unit{\m}        & 0.0068                      & 0.005                          & 0.0042                        & 0.0077                       \\
		weight / \unit{\N}          & 1.0791                      & 0.6867                         & 0.3924                        & 0.2943                       \\
		weight in water / \unit{\N} & 0.04905                     & 0.14715                        & -0.04905                      & -0.4905                      \\
		wet weight / \unit{\N}      & 1.0791                      & 0.6867                         & 0.3924                        & 0.981                        \\ \midrule[1pt]
		cable \#                    & \subref{e_acmev_es_rope}    & \subref{e_acmev_es_weighted}   & \subref{e_acmev_es_chain}     & \subref{e_acmev_es_elastic}  \\ \midrule
		diameter / \unit{\m}        & 0.0055                      & 0.0035                         & {---}                         & 0.0054                       \\
		weight / \unit{\N}          & 1.0791                      & 1.3734                         & 3.11958                       & 0.5886                       \\
		weight in water / \unit{\N} & 0.24525                     & 1.1772                         & 2.70756                       & 0.0981                       \\
		wet weight / \unit{\N}      & 1.5696                      & 1.52055                        & 3.11958                       & 0.6867                       \\ \bottomrule
	\end{tabular}
	\label{e_acmev_es_cables_characteristics}
\end{table}

\subsection{Trajectories of recorded sequences}
\label{e_acmev_trajectories_of_recorded_sequences}

All of the recorded sequences for the cables shown in \cref{e_acmev_es_all_cables} are shown in
\cref{e_acmev_tors_trajectory_6,af_rtdtec_trajectory_1,af_rtdtec_trajectory_2,af_rtdtec_trajectory_3,af_rtdtec_trajectory_4,af_rtdtec_trajectory_5,af_rtdtec_trajectory_7,af_rtdtec_trajectory_8}.
It should be noted that sequences were captured over several days, and the motion capture system
was calibrated before each session. This resulted in the world frame differing from day to day.

\begin{figure}
	\centering
	\inputtikz{figures/catenary_experiment_trajectories_6}
	\caption{Trajectories of all sequences for cable \subref{e_acmev_es_weighted}.}
	\label{e_acmev_tors_trajectory_6}
\end{figure}

The open loop characteristic of the robot's motion is clearly visible in
\cref{e_acmev_tors_trajectory_6}: the trajectories are relatively straight but are not pure surge
and sway motion with relation to the fixed point.

\section{Robot chain control experiment}
\label{e_robot_chain_control_experiment}

\subsection{Closing the control loop}
\label{e_rcce_closing_the_control_loop}

A prerequisite of any experimental application of a control system is the availability of a
measurement of the state of the system. In the simulations presented in \cref{robot_chain_control},
the measurement was assumed to be perfect, but in practice, it is not. Since this work is focused
on the control of a robot chain, it was decided to use the already available and supremely accurate
Qualisys underwater motion capture system used in the previous experiment
(\cref{e_augmented_catenary_model_experimental_validation}) and its real-time
capabilities\footnote{\url{https://docs.qualisys.com/qtm-rt-protocol/}}.

From a technical point of view, the cameras operate in a network with a computer running
specialized software - \ac{qtm} - for recording. Additionally, \ac{qtm} exposes a server which can
be queried by any other program locally or on a second computer to obtain scene data in real time;
this data can be two dimensional position of a marker in camera frames, three dimensional position
of a marker in the world frame, or even a six dimensional pose of a group of markers called rigid
body.

\begin{figure}
	\centering
	\inputtikz{figures/closed_loop_diagram}
	\caption{Principle diagram of the implementation of the closed loop.}
	\label{e_rcce_ctcl_closed_loop_diagram}
\end{figure}

A \ac{ros} package was implemented in \citeurl{filliungQtm_interface2025} to obtain the pose of a
predetermined robot and publish it in a \ac{ros} topic. This topic was then used as a close to
perfect measurement of the robot's pose in the world frame in order to command Bluerovs in a closed
loop (\cref{e_rcce_ctcl_closed_loop_diagram}). The \ac{ros} node matches the \qty{100}{\Hz}
sampling rate of \ac{qtm} which is sufficient for the control of a Bluerov, which has a relatively
slow dynamics. \marginnote{Shannon can rest easy, we are sampling fast enough.}

Additionally, a simulation twin of the node was implemented with identical interface with \ac{ros}
in order to facilitate testing of control systems relying on the Qualisys system.

\subsection{Controller implementation}
\label{ex_rcce_controller_implementation}

The \ac{mpc} controller discussed in \cref{robot_chain_control} is interfaced with \ac{ros} in
\citeurl{filliungBluerov2025}, notably a function is written to convert the force actuation
computed by the \ac{mpc} to the Bluerov's expected \ac{pwm} actuation using the robot's theoretical
maximum thrust force.

The first experimental implementation is of a single Bluerov with a pose target
(\cref{e_rcce_ctcl_mpc_diagram}). This is a simple demonstrator of the interface and was used to
pathfind the chain of \acp{rov} implementation as well as working parameters for the controller.

\begin{figure}
	\centering
	\inputtikz{figures/mpc_diagram}
	\caption{Software schematic of the single \ac{mpc} Bluerov experiment.}
	\label{e_rcce_ctcl_mpc_diagram}
\end{figure}

The second experimental implementation is of a follower Bluerov reacting to a leader Bluerov and
constraints related to the whole chain, mainly the cable between leader and follower and the pool's
dimensions. In the two Bluerov implementation, the leader's trajectory still needs to be predicted
even though it is controlled by a human in order to compute constraints on the follower's
trajectory. This is done with the pose data available in the same way as for the follower and the
leader's control input broadcasted with \ac{ros} (\cref{e_rcce_ctcl_mpc_leader_follower_diagram}).

\begin{figure}
	\centering
	\inputtikz{figures/mpc_leader_follower_diagram}
	\caption{Software schematic of the leader-follower \ac{mpc} Bluerov experiment.}
	\label{e_rcce_ctcl_mpc_leader_follower_diagram}
\end{figure}

In both of these implementations, a recovery procedure is implemented in case the \ac{mpc} fails to
converge to a solution: the tolerance of the optimization algorithm is increased until either the
problem converges and the tolerance is brought down again or the tolerance reaches a maximum value
in which case the robot as well as the experiment are stopped.

\subsection{Experiment setup}
\label{ex_rcce_experiment_setup}

The experiment takes place in the same tank as described in \cref{e_acmev_experimental_setup}
(\cref{e_acmev_es_robot_in_tank}). A larger 8-cameras Qualisys setup was used in order to cover a
larger volume than previously. The two computers running \ac{qtm} and \ac{ros} are connected via a
wireless network as opposed to a wired network because the Ethernet port of the \ac{qtm} computer
is taken by the camera network. The robots used are two Bluerovs in the Heavy configuration
\ref{e_rcce_es_rovs}.

\begin{figure}
	\centering
	\begin{subfigure}{.48\figurewidth}
		\centering
		\includegraphics[width=\textwidth, trim=900 600 2300 600, clip]{mpc_experiment_rovs.jpg}
		\caption{Basile, the leader Bluerov.}
	\end{subfigure}
	\begin{subfigure}{.48\figurewidth}
		\centering
		\includegraphics[width=\textwidth]{mpc_experiment_betty.jpg}
		\caption{Betty, the follower Bluerov.}
	\end{subfigure}
	\caption{Bluerovs used in the experiment.}
	\label{e_rcce_es_rovs}
\end{figure}

\subsubsection{One robot experiment}

In this experiment, the rate of the \ac{ros} node is set to \qty{10}{\Hz} with an horizon of
\qty{0.5}{\s}. The model of the Bluerov is set to an underactuated configuration with the pitch and
roll components of the actuation not considered in the \ac{mpc}. Furthermore, the actuation is
constant over the horizon.

The robot's pose is weighted at \num{7.5} for the position and \num{10.0} for the yaw angle. The
final pose over the horizon is additionally weighted at \num{3.0}. Additionally, a objective
function is also added to minimize the speed of the robot with a weight of \num{1.0}. The robot's
actuation is weighted at \num{2e-4} with the exception of the vertical (depth) actuation which is
weighted at \num{1e-4}.

The actuation computed by the \ac{mpc} is bound to half the maximum theoretical actuation force
that the robot is capable of. The pose of the robot is constrained to pool's size and the camera's
captation volume and the linear speed of the robot is further constrained to \qty{0.5}{\m\per\s} on
top of the objective function.

Finally the optimization problem is solved with a tolerance of \num{1e-6} and a maximum number of
iteration of \num{100}.

\subsubsection{Two robots experiment}

In this experiment the cable of the two robots chain system is a \qty{1.5}{\meter} weighted cable
(\cref{e_acmev_es_all_cables}.\subref{e_acmev_es_weighted}) and at first only considered as a
constraint and not physically attached to the robots.

Since several Bluerovs are used, the set of visual markers attached to a robot and used to compute
its pose must be unique with respect to the others (\cref{e_rcce_es_rovs}). This is achieved on a
trial and error basis.

The rate of the \ac{ros} node is again set to \qty{10}{\Hz} with an horizon of \qty{0.5}{\s}. The
model of the Bluerovs are set to an underactuated configuration with the pitch and roll components
of the actuation not considered in the \ac{mpc}. Furthermore, the actuation is constant over the
horizon. After some trial, the robot's models are fine tuned in buoyancy to match the actual
Bluerovs used.

The follower robot's pose is weighted at \num{0.0} because a specific pose is of no interest, only
the relative position with the leader. Additionally, a objective function is also added to minimize
the speed of the follower this time with a weight of \num{10.0}. The actuation is weighted at
\num{1e-6} with the exception of the yaw actuation which is weighted at \num{1.0} in order to limit
unwanted yaw motion.

The actuation computed by the \ac{mpc} is bound to half the maximum theoretical actuation force
that the robot is capable of. The pose of the robot is constrained to pool's size and the camera's
captation volume and the linear speed of the robot is further constrained to \qty{0.5}{\m\per\s} on
top of the objective function and most importantly, the inter-robot distance is constrained by the
length of the cable (\qty{1.5}{\meter}) and the latter's lowest point is constrained again by the
pool's size.

Finally the optimization problem is solved -as before- with a tolerance of \num{1e-6} and a maximum
number of iteration of \num{100}.

\subsection{Complications due to experimental conditions}
\label{ex_rcce_complications_due_to_experimental_conditions}

% \cite{mooreGeneralizedExtendedKalman2016}

\section{Conclusion}
\label{e_conclusion}